
\endnotecommand{\endnote}[1]{\textcolor{blue}{PLACEHOLDER: \textbf{#1}}}
\endnotecommand{\remove}[1]{\textcolor{red}{TO REMOVE: #1}} % For chemical formulas
\endnotecommand{\endnote}[1]{\textcolor{green}{#1}} % For chemical formulas
% \endnotecommand{\sidenote}[1]{%
%         \refstepcounter{sidenote}\mbox{\textsuperscript{\alph{sidenote}}}%
%         \marginpar{\footnotesize \mbox{\textsuperscript{\alph{sidenote}} }\textcolor{orange}{#1}}%
% }
% \endnotecounter{sidenote}
%\endnotecommand{\endnote}[1]{\textcolor{orange}{\\COMMENTS: #1}}


\endnote{Check the endnotes in BLUE. They should be replaced by a paragraph (or they are already indicating the content of an existing/subsequent paragraph)}

\remove{Remove the text in RED before submission. It is just for you to identify what to remove. It is also a way to keep some material that can be reused in other places of the manuscript.} \endnotepage

\endnote{The text in GREEN is new text}